% Parent-ready proof fragment. All labels use gt:.
% Bibliography keys used: OrderedBoundary, PolygammaResonances, BilinearTails, DLMF.
% Current source paths and verification files are listed in README_gap_ten.txt.
\section{The complete gap-ten reconstruction theorem}
\label{gt:section}

The latest incoming report \emph{Polygamma Moments and Stieltjes
Resonances} proves the complete gap-eight family and gives a computable
eventual threshold at every fixed gap \cite{PolygammaResonances}.
Thus the next unresolved fixed-gap family in the inspected material is
the gap of ten. We settle that entire family. The analytic improvement
is to keep different error budgets for odd and even Gamma modes; only
two of the five modes require the larger error. The finite proof uses
exact fast polynomial arithmetic, so a substantially longer interval
can be checked without quadratic propagation of the fundamental array.

We also give a bounded-size formula for every inverse coefficient at
an arbitrary fixed gap. At gap ten it uses determinants of order six,
independently of the two tail orders. This supplies a constructive
answer to one version of the earlier compact-inverse-row question
\cite{BilinearTails}. It retains a finite causal recurrence and does
not assert a short product evaluation of those coefficients.
These advances are relative to the inspected project sources; no
claim of historical priority is made.

\subsection{The original harmonic-tail coordinates}

Let $b\ge6$, $p=2b-10$, $q=2b$, and
\[
 A_v(n)=\zeta(v,n+1)=\zeta(v)-H_n^{(v)},\qquad x_n=n+\tfrac12.
\]
The series
\[
 D_{p,q}(s)=\sum_{n\ge0} A_p(n)A_q(n)x_n^{-s}
\]
converges absolutely for $\operatorname{Re}s>3-p-q$ and has a
meromorphic continuation. Write $D_r=D_{p,q}(-2r)$.
All these values are regular: the centered Hurwitz expansion
\begin{equation}\label{gt:tail-expansion}
 A_v(n)\sim\sum_{h\ge0}c_{v,h}x_n^{1-v-2h},\qquad
 c_{v,0}=\frac1{v-1},\quad
 c_{v,h}=\frac{B_{2h}(1/2)}{(2h)!}(v)_{2h-1}\quad(h\ge1)
\end{equation}
allows subtraction of successive terms, and the resulting possible
poles of $D_{p,q}$ are the odd integers $3-p-q-2h$.
The Bernoulli convention is $B_1=-1/2$, and $(v)_k$ is a rising factorial.

The inherited block transform is
\begin{align}
 C_j&=\sum_{r=0}^{j}\binom{2j+1}{2r}2^{-2(j-r)}D_r,
       \label{gt:C}\\
 D_r&=\sum_{j=0}^{r}\frac{\binom{2r}{2j}}{2j+1}
                      B_{2r-2j}(1/2)C_j.\label{gt:D}
\end{align}
For $j\ge b$ it has the finite late-coordinate form
\begin{equation}\label{gt:late}
 C_j=\sum_{k=0}^{b-1}a_{jk}Z_k,\qquad
 Z_k=\zeta(2k),\qquad Z_0=-\tfrac12.
\end{equation}
For $d=2j+1$ and $e=d-p-q+2k+1$, put $a_{jk}=0$ if $e<0$;
otherwise the exact coefficient is
\begin{equation}\label{gt:source-entry}
 a_{jk}=B_e\left\{
 \frac{\binom{d-q+1}{p-2k}}{d-q+1}
 +\frac{\binom{d-p+1}{q-2k}}{d-p+1}\right\}.
\end{equation}
Each summand is omitted unless its lower binomial index is between
$2$ and its upper index. In particular, a zero lower index is omitted.
This is the exact convention in the original source, and it will be
retained throughout the proof.

For context, the finite summation underlying these formulas follows
from
\[
 A_p(n-1)A_q(n-1)-A_p(n)A_q(n)
 =n^{-p}A_q(n)+n^{-q}A_p(n)+n^{-p-q}.
\]
Multiply by $n^{2j+1}$ and sum by parts. The constant term at
$N\to\infty$ of
\begin{align*}
 &N^{2j+1}A_p(N)A_q(N)
 +\sum_{n=1}^{N} n^{2j+1-q}A_p(n)\\
 &\hspace{20mm}+\sum_{n=1}^{N}n^{2j+1-p}A_q(n)
 +\sum_{n=1}^{N}n^{2j+1-p-q}
\end{align*}
is $C_j$. Expanding the centered difference
$(x+1/2)^{2j+1}-(x-1/2)^{2j+1}$ gives \eqref{gt:C};
the inverse generating factor is $z/(2\sinh(z/2))$, giving
\eqref{gt:D}. Finite Faulhaber summation gives
\eqref{gt:source-entry}. These coordinate identities and their
constant-term interpretation are inherited from the ordered-boundary
development \cite{OrderedBoundary}.

Define the original $b$-square matrix
\[
 M_{b,5}=(a_{b+r,k})_{0\le r,k<b}.
\]
On the specified formal labels
\begin{equation}\label{gt:formal-labels}
 C_0,\ldots,C_{b-1},Z_0,\ldots,Z_{b-1},
\end{equation}
the first $2b$ block rows have coefficient matrix
$\operatorname{diag}(I_b,M_{b,5})$.
The diagonal of \eqref{gt:D} is $1/(2r+1)$, so the corresponding
coefficient determinant of $D_0,\ldots,D_{2b-1}$ is
\begin{equation}\label{gt:prefix-det}
                       \frac{\det M_{b,5}}{(4b-1)!!}.
\end{equation}
Formal rank below always means rational matrix rank on
\eqref{gt:formal-labels}. It does not assert independence after
these labels are evaluated as numerical constants.

\begin{theorem}[Every gap-ten pair has the optimal prefix]
\label{gt:main}
For every integer $b\ge6$, the matrix $M_{b,5}$ is nonsingular.
Consequently, for every even $p\ge2$ and $q=p+10$, the first $q$
centered moments
\[
 D_{p,q}(0),D_{p,q}(-2),\ldots,D_{p,q}(-2(q-1))
\]
have formal rank $q$ and generate all later centered even moments
over $\mathbb Q$. The number $q$ is optimal for the $q$ specified
formal coordinates. Their coefficient determinant is
\eqref{gt:prefix-det}.
\end{theorem}

\subsection{A rational recurrence of fixed dimension}

Put
\begin{align*}
 \beta_n&=\frac{2|B_{2n}|}{(2n)!}\quad(n\ge1),\\
 a_u&=0\quad(1\le u\le5),\qquad
 a_u=(2u)^{\underline{10}}\quad(u>5),\\
 b_s&=(2s-1)^{\underline{10}},\qquad
 d_s=2(a_s+b_s)\quad(s>5).
\end{align*}
The underline denotes a falling factorial; the symbols $b_s$ here
are recurrence factors, not Bernoulli numbers. Every $d_s$ is positive.
Define the rational array
\begin{align}
 R_{u,j}&=\delta_{uj}&&(1\le u,j\le5),\notag\\
 R_{s,j}&=\frac1{d_s}\sum_{u=1}^{s-1}
             \beta_{s-u}(b_s+a_u)R_{u,j}&&(s>5),
 \label{gt:recurrence}
\end{align}
and let
\[
 F_{b,5}=(R_{b+i,j})_{1\le i,j\le5},\qquad f_{b,5}=\det F_{b,5}.
\]
The inherited complementary-minor identity is
\begin{equation}\label{gt:bridge}
 \det M_{b,5}=
 \frac{(-1)^{b(b-1)/2}}{4^bb!}
            \left(\prod_{s=6}^{b+5}d_s\right)f_{b,5}.
\end{equation}
Its normalization is important. For an even series
$R(z)=\sum_{n\ge1}r_nz^{2n}$, the coefficient of $z^{2s-11}$ in
\begin{equation}\label{gt:null}
 \mathcal T_5R=(R^{(10)}-R^{(10)}(0))\cot(z/2)
                               +(R\cot(z/2))^{(10)}
\end{equation}
is
\[
 d_sr_s-\sum_{u<s}\beta_{s-u}(b_s+a_u)r_u.
\]
The assignment $a_5=0$ records the subtracted derivative, equivalently
the omitted lower binomial index zero in \eqref{gt:source-entry}.
Truncate the five fundamental null series after degree $2b$ and
apply the complementary-minor identity to this coefficient operator.
Reversing the original $b$ columns and restoring their factorial
scales gives exactly \eqref{gt:bridge}; a full normalization proof
is in \cite[Gamma-basis section]{OrderedBoundary}.

\subsection{The inherited Gamma basis, with the normalization made explicit}

Let
\[
 W(x)=\left|\Gamma\!\left(\tfrac14+\tfrac{ix}{2}\right)\right|^{20},
 \qquad
 P_j(x)=\sinh(\pi(5-j)x)\ (1\le j<5),\qquad
 P_5(x)=\tanh(\pi x).
\]
Write
\begin{equation}\label{gt:t}
 t_{n,j}=\frac4{(2n)!}\operatorname{Im}
                  \int_0^\infty(x+i/2)^{2n}P_j(x)W(x)\,dx.
\end{equation}
These are absolutely convergent ordinary integrals. The Gamma-basis
theorem \cite{OrderedBoundary} states that
\[
 G_j(z)=4\sin(z/2)\int_0^\infty\sinh(xz)P_j(x)W(x)\,dx
\]
is a null germ of \eqref{gt:null}, and that
\begin{equation}\label{gt:gamma-bridge}
 (t_{b+i,j})_{i,j=1}^{5}=F_{b,5}T_5,
 \qquad T_5=(t_{n,j})_{n,j=1}^{5},\qquad \det T_5>0.
\end{equation}

Here is the analytic mechanism, to specify precisely what is inherited.
For an allowed multiplier $P$, set
\[
 Q_P(t)=2\sinh(t/2)\int_{\mathbb R}e^{ixt}P(x)W(x)\,dx.
\]
The Fourier contours in the two terms obtained by differentiating
$Q_P$ ten times can be shifted by $\pm i$. Gamma recurrence gives
\[
 (iy-\tfrac12)^{10}W(y+i)
 =(iy+\tfrac12)^{10}W(y-i)
 =2^{10}\left[
    \Gamma(\tfrac34+iy/2)\Gamma(\tfrac34-iy/2)\right]^{10}.
\]
The nearest poles of the common Gamma factor are at $\pm3i/2$,
outside the two shifted contours. Also $P(y+i)=P(y-i)$.
The sinh modes therefore have no contact term. For tanh the two
simple poles at $y=\pm i/2$, each with residue $1/\pi$, give
$Q_P^{(10)}(0)=i2^{12}\pi^5$; the other four top derivatives vanish.
The resulting identity is the hyperbolic version of
\eqref{gt:null}. Rotation $G_j(z)=iQ_{P_j}(iz)$ gives the stated
null equation. Stirling decay justifies all contour shifts and
differentiations: on the vertical boundaries the decay is at least
$e^{-\pi|\operatorname{Re}y|}$ times a fixed polynomial.

For the sign, $\operatorname{Im}(x+i/2)^{2n}$ has leading term
$nx^{2n-1}$. A triangular row operation and continuous Cauchy--Binet
express $\det T_5$ as a positive constant times
\[
 \int_{(0,\infty)^5}
 \Delta(x_1^2,\ldots,x_5^2)
 \Delta(\cosh\pi x_1,\ldots,\cosh\pi x_5)
 \prod_{r=1}^{5}x_rW(x_r)\tanh(\pi x_r)\,d\mathbf x.
\]
Indeed $\sinh(ku)/\tanh u=\cosh u\,U_{k-1}(\cosh u)$,
so the multiplier alternant has column degrees $4,3,2,1,0$;
their reversal has positive sign. Both Vandermondes change sign
together under any interchange, and their product is strictly
positive off the diagonals. This proves $\det T_5>0$.
The coefficient identity in \eqref{gt:gamma-bridge} then follows
from uniqueness in \eqref{gt:recurrence}.

\subsection{An effective phase estimate at half-gap five}

Put
\[
 \lambda_j=j\pi,\qquad A=2^5(2\pi)^{10},\qquad
 a_j^*=A/2\ (j<5),\qquad a_5^*=A.
\]

\begin{lemma}[Specialization of the effective Gamma estimate]
\label{gt:phase}
For every integer $N\ge147456$ and $1\le j\le5$,
\begin{align}
 \left|\frac1{a_j^*}\int_0^\infty(x+i/2)^NP_j(x)W(x)\,dx
 -e^{i\lambda_j/2}\left(
 \frac{\Gamma(N-4)}{\lambda_j^{N-4}}
 +\frac{5i}{2}\frac{\Gamma(N-5)}{\lambda_j^{N-5}}
                         \right)\right|
 <33\frac{\Gamma(N-6)}{\lambda_j^{N-6}}.\label{gt:phase-bound}
\end{align}
\end{lemma}

\begin{proof}
This is the $m=5$ case of the effective phase lemma of
\cite{PolygammaResonances}; a specialized proof follows.
For $x\ge10$, complex Stirling with its remainder \cite[\S5.11(ii)]{DLMF}
gives
\begin{equation}\label{gt:density}
 \frac{P_j(x)W(x)}{a_j^*}
 =x^{-5}e^{-\lambda_jx}(1+\delta_j(x)),\qquad
                      |\delta_j(x)|\le10x^{-2}.
\end{equation}
To check the constants, for $z=1/4+ix/2$ the remainder after
$1/(12z)$ has magnitude at most $1/(90|z|^3)$. With $t=1/(2x)$,
the logarithm of $W/(Ax^{-5}e^{-5\pi x})$ is
\[
 -\frac52\log(1+t^2)
 +5\left(\frac{\arctan t}{t}-1\right)
 +\frac{20}{3}\frac{t^2}{1+t^2}+20\operatorname{Re}\rho(z).
\]
Its absolute value is at most
$65/(24x^2)+16/(9x^3)<15/(4x^2)$.
Using $|e^u-1|\le2|u|$ in this range, and the bound
$2e^{-2\pi x}\le1/(9x^2)$ for the normalized multiplier, proves
\eqref{gt:density}. On $0\le x\le10$ one may use the generous
bound $P_j(x)W(x)/a_j^*<2^{460}$, obtained from
$\Gamma(1/4)<5$, $\pi<4$, and $e<4$.

Write $S_2=\Gamma(N-6)/\lambda_j^{N-6}$.
The density error on $x\ge\sqrt N$ is at most $20S_2$ because
\[
 |x+i/2|^N\le x^N\exp(N/(8x^2))<2x^N.
\]
On $10\le x\le\sqrt N$ its bound is
$\sqrt N(N+1)^{N/2}$.
For the model integral on $[10,\infty)$, put $y=x+i/2$ and
shift the ray down to the positive real axis. No pole is crossed.
The algebraic factor becomes $(1-i/(2y))^{-5}$, and Taylor's
theorem, for $y\ge10$, gives
\[
 \left|(1-i/(2y))^{-5}-1-\frac{5i}{2y}\right|
                         \le\frac{45}{4y^2}.
\]
For example the positive Taylor remainder is at most
$15u^2(1-u)^{-7}<45u^2$ for $u=1/(2y)\le1/20$.
Thus this part of the model error is at most $(45/4)S_2$.

The original compact integral, the short vertical segment, and the
two retained model integrals from $0$ to $10$ are together bounded
by $2^{720}15^N$. To absorb the remaining terms, set $k=N-7$.
Since $\lambda_j<20$ and $k!\ge(k/3)^k$,
\[
 S_2>\frac1{20}\left(\frac{k}{60}\right)^k.
\]
For $N\ge147456$, $k\ge3N/4$ and
$\sqrt{N+1}\le(5/4)\sqrt N$, so
\begin{align*}
 \frac{\sqrt N(N+1)^{N/2}}{S_2}
     &\le N^9(100/\sqrt N)^N<N^9 2^{-N}\le\tfrac14,\\
 \frac{2^{720}15^N}{S_2}
     &\le2^{720}N^8(1200/N)^N
                  <2^{720}N^8 4^{-N}<\tfrac14.
\end{align*}
Here $N^9\le2^{N-2}$ holds at $147456$ and persists by induction.
The total error is less than
$(20+45/4+1/2)S_2<33S_2$, proving \eqref{gt:phase-bound}.
\end{proof}

\subsection{Different determinant budgets for odd and even modes}

\begin{proposition}[A positive tail beginning at an explicit index]
\label{gt:positive-tail}
For every integer $b\ge73728$, one has $f_{b,5}>0$.
\end{proposition}

\begin{proof}
Apply Lemma~\ref{gt:phase} at $N_i=2b+2i$, $1\le i\le5$.
An odd column has leading imaginary phase $\pm1$ and relative
error bounded by
\[
 \frac{33\lambda_j^2}{(N_i-5)(N_i-6)}
                       <\frac{8448}{(N_i-6)^2},
\]
since $5\pi<16$. An even column has zero leading imaginary phase.
Its first nonzero term is the second term in
\eqref{gt:phase-bound}; its relative error is at most
\[
 \frac{66\lambda_j}{5(N_i-6)}
                       <\frac{858}{5(N_i-6)},
\]
because the largest even mode is $4\pi<13$.

Extract the positive row factors
\[
             \frac{4\Gamma(N_i-4)}{\Gamma(N_i+1)}\pi^{-2i}
\]
and positive column factors $a_j^*\lambda_j^{-2b+4}$.
In an even column extract additionally
$5\lambda_j/(2(2b-3))$. The five extracted signs are
$(+,-,-,+,+)$, with positive product. The remaining matrix is
$E_{ij}=j^{-2i}(1+\eta_{ij})$.
For its three odd columns,
\[
 |\eta_{ij}|<\frac{8448}{(2b-4)^2}<2^{-20}.
\]
For the two even columns the additional row factor is
$\rho_i=(2b-3)/(2b+2i-5)$; hence
\[
 |\eta_{ij}|
 <\frac8{2b-3}+\frac{858}{5(2b-4)}
 <\frac{898}{5(2b-4)}<2^{-9}.
\]
Both final inequalities hold already at $b=73728$ and improve
monotonically thereafter.

For $V=(j^{-2i})_{i,j=1}^{5}$, exact rational arithmetic gives
\begin{equation}\label{gt:V}
 \det V=\frac{49}{2764800000000},\qquad
 \operatorname{perm}V=\frac{106852831033}{25798901760000000000},
 \qquad
 \frac{\det V}{\operatorname{perm}V}
       =\frac{457228800}{106852831033}.
\end{equation}
Every permutation product uses exactly two even columns and three
odd columns. Therefore
\begin{align*}
 |\det E-\det V|
 &<\operatorname{perm}V
       \left\{(1+2^{-9})^2(1+2^{-20})^3-1\right\}\\
 &<\frac{\operatorname{perm}V}{240}<\det V.
\end{align*}
The two final inequalities are exact rational inequalities.
In particular $\det E>0$. The extracted row and column factors,
the product of the extracted signs, and $\det T_5$ are positive.
Equation~\eqref{gt:gamma-bridge} proves $f_{b,5}>0$.
\end{proof}

The distinction between the two error budgets is substantive: using
the worst entrywise error five times gives a weaker threshold.
The general estimate in the predecessor was $B(5)=635656$;
the preceding argument reduces the required finite interval to
$b<73728$, while using the same inherited phase remainder.

\subsection{An exact certificate for every remaining integer}

The cotangent coefficients satisfy
\begin{equation}\label{gt:beta-recurrence}
 \beta_1=\frac16,\qquad
 2(2n+1)\beta_n=\sum_{k=1}^{n-1}\beta_k\beta_{n-k}
                                      \quad(n\ge2).
\end{equation}
This follows from the Riccati equation for $(z/2)\cot(z/2)$.
We use the prime
\[
                         \ell=998244353=119\cdot2^{23}+1.
\]
The supplied verifier proves primality by trial division. It also
verifies that $3$ is a primitive root by checking the three prime
divisors $2,7,17$ of $\ell-1$. Every polynomial transform used in
the certificate has length at most $2^{18}$, so all required roots
of unity exist in $\mathbb F_\ell$.

For $6\le b\le73727$, the largest needed fundamental row is
$73732$ and the largest needed cotangent index is $73731$.
All recurrence denominators are invertible modulo $\ell$.
This is immediate for $2(2n+1)$ in \eqref{gt:beta-recurrence};
for the other diagonal use
\begin{equation}\label{gt:d-factor}
 d_s=2\frac{(2s)!}{(2s-10)!}\frac{2s-5}{s},\qquad s\ge6.
\end{equation}
Every displayed integer factor is positive and less than $\ell$.
Induction therefore identifies the computed residues with the
reductions of the exact rational coefficients.

\begin{proposition}[Finite gap-ten certificate]
\label{gt:finite}
For every integer $6\le b\le73727$, the rational determinant
$f_{b,5}$ has a nonzero reduction modulo $998244353$.
\end{proposition}

\begin{proof}[Exact computational proof]
The program \texttt{verify\_gap\_ten.cpp} constructs the necessary
cotangent coefficients and five fundamental columns using integer
arithmetic only. It then evaluates all $73722$ five-by-five
determinants by Gaussian elimination over $\mathbb F_\ell$.
Every residue is nonzero. The full indexed list is
\texttt{gap\_ten\_certificate.csv}.
\begin{samepage}
\noindent Its SHA-256 identifier is
\begin{center}\small\ttfamily
f3f749a4c4e02c27b32a58e1f65714bab568a08f3b0ebdb0fd8cb8d4bf74184d
\end{center}
\end{samepage}
The recurrence and the nonzero residues are the certificate; its
hash serves only to identify the supplied record.

Every fundamental recurrence is independently checked after the
causal computation by two full polynomial products per column.
The cotangent coefficients are independently checked against
\eqref{gt:beta-recurrence}, after their initial construction by
formal division of cosine by normalized sine.
The separate Python verifier, which does not call the NTT code,
recomputes $506$ boundary determinants with the direct quadratic
recurrence, checks $64$ cotangent coefficients against rational
Bernoulli numbers, and compares the original Bernoulli matrix
\eqref{gt:source-entry} and its full normalization
\eqref{gt:bridge} for every $6\le b\le18$.
All these checks pass; their exact values are supplied in
\texttt{gap\_ten\_independent\_checks.json}.
\end{proof}

Propositions~\ref{gt:positive-tail} and \ref{gt:finite}, followed
by \eqref{gt:bridge} and \eqref{gt:prefix-det}, prove
Theorem~\ref{gt:main}. The finite-field part proves nonvanishing,
not the real signs of the corresponding rational numbers.
Indeed direct rational checks give $f_{6,5}>0$, $f_{7,5}<0$,
and $f_{18,5}>0$. A blanket positivity claim at every admissible
index would be false.

\subsection{Why the finite recurrence can be verified efficiently}
\label{gt:fast}

The accelerated calculation has a short exact algebraic description.
Define formal series
\[
 S(z)=\sum_{n\ge0}\frac{(-1)^n z^n}{4^n(2n+1)!},\qquad
 C(z)=\sum_{n\ge0}\frac{(-1)^n z^n}{4^n(2n)!}.
\]
Since $S(0)=1$, its inverse exists over $\mathbb F_\ell$, and
\[
             1-\tfrac12\sum_{n\ge1}\beta_nz^n=C(z)/S(z).
\]
Newton iteration $U\mapsto U(2-SU)$ doubles the number of correct
inverse coefficients. The number-theoretic transform computes each
polynomial product exactly in the finite field; it introduces no
roundoff error.

For a fundamental column $r_u=R_{u,j}$, the recurrence is
\[
 d_sr_s=b_s\sum_{u<s}\beta_{s-u}r_u
                  +\sum_{u<s}\beta_{s-u}a_ur_u.
\]
Both sums are causal convolutions. Divide an index interval into
two equal halves, compute its left half, add both left-to-right
convolutions, and then compute its right half. Induction on the
interval length proves that every already determined coefficient
contributes exactly once to every later equation. The transforms
for a given interval length use the same cotangent kernel and can
be reused across all five columns.

In a cyclic transform of length $L$, the left input is supported
in $[0,L/2)$ and the kernel in $[0,L)$. All wrapped terms land in
$[0,L/2)$; none contaminates the right half being updated.
This proves the validity of the shorter cyclic convolution used
by the implementation. For a fixed half-gap $m$, the resulting
cost is $O(mN\log^2N+Nm^3)$ field operations and $O(mN)$ stored
field elements, using ordinary cubic elimination for the consecutive
boundary determinants. The same method applies to any prescribed
half-gap; it is not tied to five columns.

\section{Inverse harmonic coordinates from bounded-size determinants}
\label{gt:green-section}

For this section let $m\ge1$ be arbitrary, $b>m$, and use the same
definitions with $10$ replaced by $2m$ and $5$ by $m$.
Thus $p=2b-2m$, $q=2b$, $F_{b,m}=(R_{b+i,j})_{i,j=1}^m$, and
assume $f_{b,m}=\det F_{b,m}\ne0$. This assumption is now a theorem
for every $b>5$ when $m=5$.

For an integer forcing index $N>m$, define the causal Green column
$\mathsf G_{s,N}$ by zero initial values for $1\le s\le m$ and
\begin{equation}\label{gt:Green}
 d_s\mathsf G_{s,N}
 =\delta_{sN}+\sum_{u<s}\beta_{s-u}(b_s+a_u)\mathsf G_{u,N}
                                  \qquad(s>m).
\end{equation}
In particular $\mathsf G_{s,N}=0$ for $s<N$ and
$\mathsf G_{N,N}=1/d_N$. Put
\[
 \mathbf r_n=(R_{n,1},\ldots,R_{n,m}),\qquad
 \mathbf g_N=(\mathsf G_{b+1,N},\ldots,
                         \mathsf G_{b+m,N})^{\mathsf T}.
\]

\begin{theorem}[Every inverse coefficient in dimension $m+1$]
\label{gt:inverse}
For $0\le k,r<b$, set $n=b-k$ and $N=r+m+1$. Then
\begin{equation}\label{gt:inverse-formula}
 \boxed{\displaystyle
 (M_{b,m}^{-1})_{k,r}
 =\frac{2(-1)^{n+N}(2n)!}{(2r+1)!}\,
 \frac{\det\begin{pmatrix}
       F_{b,m}&\mathbf g_N\\
       \mathbf r_n&\mathsf G_{n,N}
                  \end{pmatrix}}{\det F_{b,m}}.}
\end{equation}
The determinant in the numerator has order $m+1$, independent of
$b$. Every entry is rational and specified by a finite recurrence.
\end{theorem}

\begin{proof}
Let $J_{m,b}$ be the coefficient matrix of the cotangent operator
\eqref{gt:null}, with $10$ replaced by $2m$, acting on
$\sum_{n=1}^b r_nz^{2n}$ and with output rows
$N=m+1,\ldots,b+m$. Its equation with unit forcing at $N$ is
the recurrence \eqref{gt:Green}, with the boundary conditions
$r_{b+1}=\cdots=r_{b+m}=0$ imposed after truncation.
The general full solution is
\[
          r_s=\mathbf r_s\mathbf c+\mathsf G_{s,N}.
\]
Its $m$ terminal conditions give
$F_{b,m}\mathbf c=-\mathbf g_N$. Hence
\[
 (J_{m,b}^{-1})_{n-1,r}
 =\mathsf G_{n,N}-\mathbf r_nF_{b,m}^{-1}\mathbf g_N
 =\frac{\det\begin{pmatrix}F_{b,m}&\mathbf g_N\\
                  \mathbf r_n&\mathsf G_{n,N}\end{pmatrix}}
        {\det F_{b,m}}.
\]

To restore the original scales, let $K_{m,b}$ be the analogous
hyperbolic-cotangent coefficient matrix. If $M^{\mathrm{rev}}$
denotes column reversal, direct comparison with
\eqref{gt:source-entry} gives
\[
 K_{m,b}=\operatorname{diag}_{0\le r<b}
           \left(\frac2{(2r+1)!}\right)
         M_{b,m}^{\mathrm{rev}}
         \operatorname{diag}_{1\le n\le b}\bigl((2n)!\bigr).
\]
Since $2B_{2h}/(2h)!=(-1)^{h-1}\beta_h$ for $h\ge1$,
its entries satisfy
$(K_{m,b})_{r,n-1}=(-1)^{N+n}(J_{m,b})_{r,n-1}$.
The same row and column signs therefore occur in the inverse.
Finally the original row $k$ corresponds to the reversed index
$n=b-k$. These factors give precisely \eqref{gt:inverse-formula}.
\end{proof}

This formula constructs each inverse row without expanding a
$b$-square determinant. It does not remove the need to compute the
finite Green column, and it makes no minimality claim about the
determinant size. The independent verifier checks all $293$ entries
of the resulting inverse matrices at $b=6,7,8,12$ directly against
the original source matrices, and separately checks an actual
six-by-six bordered determinant.

\section{Ordinary polygamma identities at every gap of ten}
\label{gt:identities}

Let $b\ge6$, $p=2b-10$, $q=2b$, and let
$\lambda_{b+r}=(M_{b,5}^{-1})_{0,r}$ for $0\le r<b$.
Define
\begin{align}
 K_b(y)&=\sum_{j=b}^{2b-1}\lambda_jy^{2j+1},\notag\\
 P_b(x)&=K_b(x+\tfrac12)-K_b(x-\tfrac12)
                   =\sum_{r=0}^{q-1}w_rx^{2r}.\label{gt:P}
\end{align}
By Theorem~\ref{gt:inverse}, all its coefficients can be given by
six-by-six determinants and the five-column recurrence.
With $c_{v,h}$ from \eqref{gt:tail-expansion}, put
\[
 u_h=\sum_{e=0}^{h}c_{p,e}c_{q,h-e}\quad(0\le h\le5),
 \qquad
 S_b(x)=\sum_{h=0}^{5}
       \left(\sum_{e=0}^{h}u_e w_{q-1-h+e}\right)x^{10-2h}.
\]

\begin{corollary}[An all-order rational family]
\label{gt:ordinary}
For every integer $b\ge6$,
\begin{equation}\label{gt:ordinary-sum}
 \boxed{\displaystyle
 \sum_{n=0}^{\infty}\left\{
 \frac{P_b(n+1/2)}{(p-1)!(q-1)!}
   \psi^{(p-1)}(n+1)\psi^{(q-1)}(n+1)-S_b(n+1/2)
                         \right\}=-\frac12.}
\end{equation}
The sum is absolutely convergent. The polynomial $P_b$ has exact
degree $2q-2$; its leading coefficient is
\begin{equation}\label{gt:leading}
 [x^{2q-2}]P_b(x)=
       \frac{4b(4b-1)f_{b-1,5}}{d_{b+5}f_{b,5}}\ne0,
\end{equation}
where at the initial index one uses the auxiliary recurrence value
\begin{equation}\label{gt:auxiliary}
 f_{5,5}=
 \frac1{363698159177574318706747080134319581355520819200000000000}.
\end{equation}
\end{corollary}

\begin{proof}
The inverse row gives $\sum_{j=b}^{2b-1}\lambda_jC_j=Z_0=-1/2$.
Expanding the centered difference in \eqref{gt:P} and using
\eqref{gt:C} gives $\sum_rw_rD_r=-1/2$.
By \eqref{gt:tail-expansion}, the nonnegative powers in
$P_b(x_n)A_p(n)A_q(n)$ are exactly $S_b(x_n)$, and the remaining
summand is $O(x_n^{-2})$. The continued sums of the subtraction
monomials vanish because $\zeta(-2h,1/2)=0$ for every $h\ge0$.
Equivalently, near spectral parameter zero, subtracting this finite
polynomial gives a normally convergent Dirichlet series and leaves
only those Hurwitz-zeta terms. Its value at zero is the ordinary
sum of the subtracted terms. Finally
$A_v(n)=\psi^{(v-1)}(n+1)/(v-1)!$ for positive even $v$.

For $b\ge7$, delete the last row and first column of $M_{b,5}$.
Reindexing the two orders down by two and the two source indices
down by one identifies the minor exactly with $M_{b-1,5}$,
including the omitted-binomial convention. Cramer's rule and
\eqref{gt:bridge} give
\[
 \lambda_{2b-1}=(-1)^{b-1}\frac{\det M_{b-1,5}}{\det M_{b,5}}
               =\frac{4b f_{b-1,5}}{d_{b+5}f_{b,5}}.
\]
The centered difference multiplies this highest coefficient by
$4b-1$, proving \eqref{gt:leading}. At $b=6$ direct rational
recurrence gives \eqref{gt:auxiliary}, and the original matrix
gives $\lambda_{11}=132/127$, verifying the same formula without
introducing a zero-order tail. Nonvanishing follows from
Theorem~\ref{gt:main} and the explicit auxiliary value.
\end{proof}

\subsection[The newly covered pair (14,24)]{The newly covered pair $(14,24)$}

Take $b=12$, so $p=14$ and $q=24$. Formula \eqref{gt:ordinary-sum}
is then the explicit ordinary identity
\[
 \sum_{n\ge0}\left\{
 \frac{P_{12}(n+1/2)}{13!\,23!}
 \psi^{(13)}(n+1)\psi^{(23)}(n+1)-S_{12}(n+1/2)
                         \right\}=-\frac12,
\]
with $\deg P_{12}=46$ and $\deg S_{12}=10$. Its leading coefficients are
\begin{align*}
 [x^{46}]P_{12}(x)
 &=\frac{196667260793027525874874688960}
          {14009065684434488629739697248281},\\
 [x^{10}]S_{12}(x)
 &=\frac{196667260793027525874874688960}
          {4188710639645912100292169477236019}.
\end{align*}
The second denominator is $13\cdot23$ times the first.
The file \texttt{gap\_ten\_14\_24\_example.json} contains all $24$
even coefficients of $P_{12}$, all six coefficients of $S_{12}$,
and the twelve coefficients of $K_{12}$ as exact rationals.
The formula \eqref{gt:inverse-formula} is an alternative finite
specification of the same coefficients. The generator checks the
inverse against the original Bernoulli matrix and the entire
subtraction polynomial symbolically. This pair lies beyond the
previous smaller-order-through-twelve coverage.

More generally, replace row zero by row $k$ of the inverse in
\eqref{gt:P}. The identical construction, with the same degree
bound and subtraction rule, gives the right-hand side
$\zeta(2k)$ for every $0\le k<b$. Thus the completed reconstruction
theorem gives a family of ordinary convergent polygamma-product
representations for every even zeta coordinate in the source block.

